\section{调用模式与执行生命周期}
\subsection{发现、已知 Provider 与本地组合}
当执行者尚未确定时使用普通发现。User 在有界窗口内收集 ACK，并从通过验证与策略检查的 positive ACK 中选择。Negative ACK 表示不参与；未收到 ACK 则不能区分无意参与、丢包或断连。选择策略决定本操作的候选。在支持多个独立执行的情况下，选择多个执行者仍不同于在协作计划中分配具有依赖的互补角色。

当应用含义固定设备时使用指定 Provider 调用，例如停止之前启动的相机。当前 \texttt{RequestServiceTargeted} 路径在支持的条件下复用预先建立的一次性授权材料，避免重复 ACK/Selection，但仍检查当前授权和请求状态。快速路径并不适用于所有模式：请求级保密交付需要 Selection 携带输入和结果密钥，当前实现对此使用 bootstrap/ACK/Selection 路径。显式 Selection-gated input 模式则被不含 Selection 的 Targeted 接口拒绝。因此，指定设备属于应用语义，省略交换只是有条件的实现优化。

受信本地组合通过进程内 local registry 调用，既没有远程发现，也不具有相同网络威胁边界。评价需要记录处理函数是由本地调用、普通发现还是指定调用到达，否则看似调用性能改善可能只是测量了另一种操作。

\subsection{准入、进度与失败}
API 提交、进入 runtime 队列、被 Provider 接受及应用工作完成是不同事件。过载时只报告完成请求的延迟，会掩盖丢弃或拒绝。因此 runtime 需要可观测的拒绝、超时与终态结果，评价应在明确分母中计入所有提交。Selection 状态查询有助于了解选定工作的已观测状态，但不证明物理动作已发生，也不提供分布式事务。

应用还必须区分不再等待结果与撤销实际动作。录制可以停止，但已记录的帧仍存在；推理 worker 可以停止产生新输出，但删除请求者状态不能撤回已完成的远程动作。重试需要明确的幂等或补偿策略。拟议协作扩展中的替换还要求旧分配和晚到回调不再推进当前计划。

\section{控制消息、数据流与完整对象}
小型服务消息表达意图、权限核验材料与状态。大应用数据不应复制到每条控制消息中。对象引用使服务可以标识模型分片、录制文件或图像，而接收者独立获取对应命名 Data。引用和服务契约应提供选定版本及验证重建对象所需的信息。具体接口决定可用的名称、大小、摘要、分段和保护元数据，并非所有对象都需要完全相同字段。

标识对象、获取字节和授权使用应当分开。只有预期摘要本身可信时，摘要比对才可用于判断字节是否符合预期。签名按照应用信任规则将 Data 关联到密钥。调用或协作检查决定生产者和对象能否满足当前输入，接收者专属密钥交付限制对受保护 Content 的访问。这些检查不能相互替代。

持续流与完整对象具有不同生命周期。直播视频可以优先近期帧，并容忍明确标记的缺口；模型文件通常需要取得全部必要字节才能加载。增量推理需要有序应用事件及终态，但 token 输出流未必与相机流使用可互换接口。NDNSF 提供命名数据机制，应用定义可容忍缺口、缓冲上限和完整性规则。评价必须计入所选策略的丢弃和排队成本。

多视角场景体现命名对象复用：计算 Provider 返回标注和原始图像引用，并交付获准的图像内容密钥。请求者从可用生产者或缓存获取图像，验证原始签名，而无需仅因接收者变化就重新生成图像对象。这不保证缓存副本始终可用；长期保留需要明确的生产者或 repository 策略。
